\SourcePage{398}{401}
\section{Multiplet terms. Cases $c$ and $d$}

Besides coupling cases $a$ and $b$, and those intermediate between them, other types of coupling also exist. Their origin is as follows. The appearance of the quantum
\SourcePage{399}{402}
number $\Lambda$ is due, in the final analysis, to the electric interaction between the two atoms in the molecule, which gives axial symmetry to the problem of determining the electronic terms (this interaction in a molecule is referred to as coupling of the orbital angular momentum to the axis). The magnitude of this interaction is measured by the separations between terms having different values of $\Lambda$. In all the preceding discussion, this interaction was tacitly assumed to be so strong that these separations are large both in comparison with the intervals of the multiplet splitting and with those of the rotational structure of the terms. There are, however, also converse situations in which the coupling of the orbital angular momentum to the axis is comparable with, or even small in comparison with, other effects. In such situations, of course, the projection of the orbital angular momentum on the axis cannot be regarded as conserved in any approximation, and the number $\Lambda$ therefore loses its meaning.
